\section{Design Rationale: ADHD Assist Pillars}
\label{sec:design}

\begin{table*}[t!]
\caption{Technique $\times$ symptom mesh for ADHD Assist. Cells are design
ratings from prior evidence and our synthesis, not Study~1 clinical outcomes.
Columns: Sust.\ = sustained attention; WM = working memory; Init.\ = task
initiation; Imp.\ = impulsivity; Infl.\ = inflexibility; Rew.\ =
reward~/~motivation; Emot.\ = emotional dysregulation; Meta.\ = metacognition.}
\label{tab:mesh}
\centering
\footnotesize
\setlength{\tabcolsep}{4.5pt}
\begin{tabular}{l*{9}{c}}
\toprule
Technique & Sust. & WM & Init. & Imp. & Time & Infl. & Rew. & Emot. & Meta. \\
\midrule
P1 Concise & \textbf{S} & \textbf{S} & \textbf{S} & I & P & I & P & I & n \\
P2 Structured & \textbf{S} & \textbf{S} & \textbf{S} & P & P & P & I & I & P \\
P3 Progressive & \textbf{S} & \textbf{S} & \textbf{S} & P & P & I & P & I & P \\
P4 Single-focus & \textbf{S} & P & P & \textbf{S} & I & \textbf{S} & I & I & P \\
P5 Gentle redirect & \textbf{S} & I & n & \textbf{S} & P & \textbf{S} & I & \textbf{S} & P \\
\bottomrule
\end{tabular}
\end{table*}

If default LLM tutoring fails students with ADHD, what should a reply look
like instead?

ADHD Assist uses five \textbf{response-shape} pillars. They draw on cognitive
load and accessibility research~\cite{sweller2011cognitive,cowan2010magical,w3c2021coga},
ADHD executive-function accounts~\cite{barkley1997behavioral}, and ADHD
co-design evidence from Zhu, Yu, and Luo~\cite{zhu2026scaffolding}. We mapped
each technique to the ADHD difficulties it is meant to help. The pillars
become the Teacher policy and the Dean's written constitution
(\S\ref{sec:system}).

Study~1 tests whether the structural signature of these pillars holds under
multi-turn use when a Dean is present; it does not re-test every symptom
claim as a clinical outcome. The mesh ratings in this section
(Strong~/~Partial~/~Indirect) are likewise design ratings from the literature
and our synthesis.

\subsection{The five pillars}
\label{sec:pillars}

\paragraph{P1. Concise.}
Working memory only holds a few chunks at once~\cite{cowan2010magical}. Long
undifferentiated text raises extraneous load~\cite{sweller2011cognitive} and
is hard to re-enter after an interruption~\cite{zhu2026scaffolding}.

\emph{In practice:} tutoring aims about 150 words (hard cap 250); clarifications
shorter. No filler praise. If the topic exceeds the cap, summarize and invite
continuation (P3).

\emph{When it fails:} lectures, repeated reassurance, answer buried under
throat-clearing.

\paragraph{P2. Structured.}
Headings, bullets, and a fixed order externalize ``where am I?'' so the
learner does not rebuild context from prose after
interruption~\cite{w3c2021coga,barkley1997behavioral}. SocraticLM scores tutor
readability for the same reason~\cite{liu2024socraticlm}.

\emph{In practice:} scaffolded turns open with a short summary block and close
with one continuation invite; optional step ladder (at most five) and at most
one quick check. Scoring anchors are Top summary and Next?. The interface may
label those blocks TLDR and Continue; the underlying checks are unchanged.

\emph{When it fails:} free-form essay; missing summary or Next?; nested
sprawl.

\paragraph{P3. Progressively disclosed.}
One decision surface at a time beats dumping everything up
front~\cite{saha2023can,zhu2026scaffolding}.

\emph{In practice:} answer the immediate ask tightly first. Add depth only
after an explicit Next? invite. If the user asks two questions, answer the
first and defer the second.

\emph{When it fails:} full plan or lecture in one shot; multiple open threads.

\paragraph{P4. Single-focus.}
Multiple agendas leave attention residue and invite topic-merge when
inhibition is weak~\cite{sweller2011cognitive,zhu2026scaffolding,shani2024multiturn}.

\emph{In practice:} one main topic per turn. Steps serve that topic only.

\emph{When it fails:} two unrelated headings; silently answering a second
injected topic.

\emph{Limit:} policy and redirect handling encode this intent. Automatic
topic-count detection is weak, so Study~1 uses scripted interrupt probes
rather than relying on an automatic one-topic detector.

\paragraph{P5. On-task continuity (gentle redirect).}
Off-topic jumps are common for many students with ADHD. Harsh refusal raises
frustration~\cite{beheshti2020emotion,zhu2026scaffolding}. SocraticLM scores
productive handling of irrelevant moves~\cite{liu2024socraticlm}.

\emph{In practice:} acknowledge the jump, name the prior topic briefly, offer
return or switch, without scolding.

\emph{When it fails:} silent tangent-follow, or a blunt ``I can't help with
that.''

\subsection{Supporting constraints (not separate IVs)}

Still enforced, but not the main style attributes under test:
\begin{itemize}
  \item \textbf{Validate and move:} confirm a correct grasp, then advance
    (SocraticLM validate-and-move analogue~\cite{liu2024socraticlm}).
  \item \textbf{Honest unknowns:} do not invent course facts.
  \item \textbf{No clinical inference} from chat.
  \item \textbf{Learner-owned reflection:} optional; off in Study~1 so style
    remains the only planned difference.
\end{itemize}

The Dean (\S\ref{sec:system}) audits the pillars through a written
constitution. It is not a sixth pillar.


\subsection{Technique $\times$ symptom mesh}
\label{sec:mesh}

Prior ADHD, load, and co-design work already links reply techniques to
specific difficulties: short chunked answers to working-memory
limits~\cite{cowan2010magical,sweller2011cognitive}; external structure to
re-entry and attention~\cite{barkley1997behavioral,w3c2021coga}; progressive
disclosure to initiation and overload~\cite{saha2023can,zhu2026scaffolding};
single-focus to impulsivity and topic-merge~\cite{shani2024multiturn}; gentle
redirect to disruption and emotional
dysregulation~\cite{zhu2026scaffolding,beheshti2020emotion,liu2024socraticlm}.

We assemble those links into one technique $\times$ deficit mesh for ADHD
Assist and rate each cell from that evidence. Table~\ref{tab:mesh} is the
design map (\textbf{S}~=~strong intended help; \textbf{P}~=~partial;
\textbf{I}~=~indirect; \textbf{n}~=~not claimed). Systems that ship structure
or a second-pass checker usually motivate the technique; they less often
publish this matrix as an explicit product constitution.

P1--P3 carry the working-memory, attention, and initiation core. P4--P5 carry
impulsivity and set-shifting when the session derails. Emotional dysregulation
is addressed mainly by redirect tone (P5), not a separate affect model.
Reward~/~motivation and deep metacognition are weakly covered here; richer
engagement pedagogy sits outside this paper's independent variable.

\subsection{Hand-off to system and Study~1}

Pillars become Teacher policy. The constitution becomes the Dean checklist.
Turn profiles decide when the full scaffold applies versus a redirect.
Study~1 tests whether that structural signature sticks better with a Dean
than with prompting alone.
